\section{引言}

线性微分方程在数学的发展中占据着重要地位。传统上，人们主要利用分析方法研究其解的性质，例如解的存在性、唯一性以及渐近行为等。然而，除了分析视角之外，还可以从代数的角度理解微分方程，将其视为某种代数结构上的作用问题。这种观点为微分方程的研究提供了新的工具和思路。

一个基本的观察是，线性微分方程可以通过微分算子来表达。例如，在一元情形下，微分算子
\[
    \partial = \frac{d}{dx}
\]
作用在函数空间上。方程
\[
    (\partial - \lambda)f = 0
\]
可以理解为一个微分算子对函数的作用结果为零。这启发我们将“微分方程”视为“算子作用的代数关系”。由此，自然地引出对“微分算子所构成的代数结构”的研究。
$\mathcal{D}$-模理论正是在这一思想基础上发展起来的。



完整的 $\mathcal{D}$-模理论具体是在$n$维光滑代数簇$X$上，$\mathcal{D}_X$为$X$上的微分算子层，$\mathcal{D}$-模则是$\mathcal{D}$-模层。通过研究$\mathcal{D}$-模的结构，可以揭示线性微分方程解空间的性质，例如解的维数、单值性以及解析延拓等问题。
涉及代数几何、同调代数等较为深刻的理论工具，但其中的一些基本思想可以在较为初等的代数框架下加以理解。在仿射情形（$X=\Spec k[x_1,\ldots,x_n]$）下，微分算子环可以具体描述为 Weyl 代数$A_n$，它由乘法算子与微分算子生成，并满足一定的交换关系。对 Weyl 代数上的模进行研究，便构成了 $\mathcal{D}$-模理论的基本内容。

本文旨在从代数的角度出发，对$\mathcal{D}$-模的基本概念进行综述。我们将首先介绍相关的代数预备知识，然后构造 Weyl 代数并说明其与线性微分方程之间的联系，通过若干具体例子展示$\mathcal{D}$-模的基本结构。最后介绍 holonomic $\mathcal{D}$-模以及 Bernstein–Sato 多项式（$b$-函数）的基本思想，以说明该理论在更深入研究中的重要性。

